\section*{\texorpdfstring{\S\CurrentExerciseGroup}{Section \CurrentExerciseGroup}}
\phantomsection
\addcontentsline{toc}{section}{\texorpdfstring{Exercises for \S\CurrentExerciseGroup}{Exercises for Section \CurrentExerciseGroup}}

\begin{enumerate}
\def\labelenumi{\arabic{enumi})}
\tightlist
\item
  Let \(X, Y\) be two locally compact spaces.
\end{enumerate}

\begin{enumerate}
\def\labelenumi{\alph{enumi})}
\item
  Show that the mapping \((\mu, \nu) \mapsto \mu \otimes \nu\) of \(\mathcal{M}(\mathrm{X};\mathbf{C}) \times \mathcal{M}(\mathrm{Y};\mathbf{C})\) into \(\mathcal{M}(\mathrm{X} \times \mathrm{Y};\mathbf{C})\) is continuous when \(\mathcal{M}(\mathrm{X};\mathbf{C})\) , \(\mathcal{M}(\mathrm{Y};\mathbf{C})\) and \(\mathcal{M}(\mathrm{X} \times \mathrm{Y};\mathbf{C})\) are equipped with the quasi-strong topology (§1, Exer. 8).
\item
  Show that the mapping \((\mu, \nu) \mapsto \mu \otimes \nu\) of \(\mathcal{M}_{+}(X) \times \mathcal{M}_{+}(Y)\) into \(\mathcal{M}_{+}(X \times Y)\) is continuous when \(\mathcal{M}_{+}(X)\) , \(\mathcal{M}_{+}(Y)\) and \(\mathcal{M}_{+}(X \times Y)\) are equipped with the vague topology.
\end{enumerate}

¶ 2) a) Let X be the interval {[}0,1{]} of R; show that if \((a_k)_{1 \leqslant k \leqslant n}\) is any finite sequence of pairwise distinct elements of X, then the \(n\) functions \(|x - a_k|\) ( \(1 \leqslant k \leqslant n\) ) are linearly independent. From this, deduce that the continuous function \(|x - y|\) defined on \(X \times X\) is not of the form \(\sum_{i=1}^{n} u_i(x)v_i(y)\) , where the \(u_i\) and \(v_i\) are continuous.

\begin{enumerate}
\def\labelenumi{\alph{enumi})}
\setcounter{enumi}{1}
\item
  For every measure \(\mu \in \mathcal{M}(\mathrm{X};\mathbf{C})\) , set \(\mathbf{f}_{\mu}(y) = \int |x - y| d\mu(x)\) . Show that the linear mapping \(\mu \mapsto \mathbf{f}_{\mu}\) of \(\mathcal{M}(\mathrm{X};\mathbf{C})\) into \(\mathcal{C}(\mathrm{X};\mathbf{C})\) is injective, has for image a dense linear subspace L of \(\mathcal{C}(\mathrm{X};\mathbf{C})\) , and is continuous but not bicontinuous for the normed space topologies on \(\mathcal{M}(\mathrm{X};\mathbf{C})\) and L. (Observe that L contains the constants and the piecewise linear functions.)
\item
  Let \(u_{i}\) ( \(1 \leqslant i \leqslant m\) ) be any \(m\) functions in \(\mathcal{C}(\mathrm{X};\mathbf{C})\) and let \(V\) be the linear subspace of \(\mathcal{M}(\mathrm{X};\mathbf{C})\) orthogonal to the \(u_{i}\) , which is therefore of codimension \(\leqslant m\) in \(\mathcal{M}(\mathrm{X};\mathbf{C})\) . Show that if \(B\) is the set of measures \(\mu\) on \(X\) such that \(\| \mu \| \leqslant 1\) , there exist measures \(\mu \in B\) and \(\nu \in V\) such that \(\int \mathfrak{f}_{\mu}(y)d\nu(y)\) is arbitrarily large (make use of \(b\) ). From this, deduce that if \(B\) , \(\mathcal{M}(X; \mathbf{C})\) and \(\mathcal{M}(X \times X; \mathbf{C})\) are equipped with the vague topology, then the mapping \((\mu, \nu) \mapsto \mu \otimes \nu\) of \(B \times \mathcal{M}(X; \mathbf{C})\) into \(\mathcal{M}(X \times X; \mathbf{C})\) is not continuous.
\end{enumerate}

¶ 3) a) Let X, Y be two locally compact spaces, K ⊂ X, L ⊂ Y two compact sets, A a compact subset of the space \(\mathcal{K}(X \times Y, K \times L; \mathbf{C})\) . For every integer \(n\) , let \(V_n\) be an entourage of the uniform structure of K such that, for every pair \((x', x'') \in V_n\) , every \(y \in L\) and every function \(f \in A\) ,

\[
\left| f (x ^ {\prime}, y) - f (x ^ {\prime \prime}, y) \right| \leqslant 1 / n ^ {2}.
\]

Let \(C_n\) be the set of functions \(y \mapsto n(f(x', y) - f(x'', y))\) for all pairs \((x', x'') \in V_n\) and every \(f \in A\) ; show that the union \(C\) of the \(C_n\) for all \(n \geqslant 1\) is a relatively compact subset of \(\mathcal{K}(Y, L; \mathbf{C})\) .

\begin{enumerate}
\def\labelenumi{\alph{enumi})}
\setcounter{enumi}{1}
\tightlist
\item
  Let B be the union of C and the set of mappings \(y \mapsto f(x, y)\) for \(x \in K\) and \(f \in A\) , which is a relatively compact subset of \(\mathcal{K}(\mathrm{Y}, \mathrm{L}; \mathbf{C})\) . Show that when \(\nu\) runs over the polar \(B^{\circ}\) of B in \(\mathcal{M}(\mathrm{Y}; \mathbf{C})\) and f runs over A, the set of mappings
\end{enumerate}

\[
y \mapsto \int f (x, y) d \nu (y)
\]

is relatively compact in \(\mathcal{H}(\mathrm{X},\mathrm{K};\mathbf{C})\)

\begin{enumerate}
\def\labelenumi{\alph{enumi})}
\setcounter{enumi}{2}
\tightlist
\item
  Deduce from b) that the mapping \((\mu, \nu) \mapsto \mu \otimes \nu\) of \(\mathcal{M}(\mathrm{X}; \mathbf{C}) \times \mathcal{M}(\mathrm{Y}; \mathbf{C})\) into \(\mathcal{M}(\mathrm{X} \times \mathrm{Y}; \mathbf{C})\) is continuous when \(\mathcal{M}(\mathrm{X}; \mathbf{C})\) , \(\mathcal{M}(\mathrm{Y}; \mathbf{C})\) and \(\mathcal{M}(\mathrm{X} \times \mathrm{Y}; \mathbf{C})\) are equipped with the topology of strictly compact convergence.
\end{enumerate}

\begin{enumerate}
\def\labelenumi{\arabic{enumi})}
\setcounter{enumi}{3}
\item
  Let X be a locally compact space, Y a paracompact locally compact space. Show that the canonical mapping of \(\mathcal{K}(X \times Y; \mathbf{C})\) into \(\mathcal{K}(X; \mathcal{K}(Y; \mathbf{C}))\) is an isomorphism of topological vector spaces.
\item
  Let \((\mathrm{X}_{\lambda})_{\lambda \in L}\) be an infinite family of compact spaces; for each \(\lambda \in L\) , let \(a_{\lambda}\) be a point of \(\mathrm{X}_{\lambda}\) and let \(\mu_{\lambda}\) be a positive measure on \(\mathrm{X}_{\lambda}\) of total mass 1. For every finite subset \(J\) of \(L\) , let \(\nu(J)\) be the measure on \(X = \prod_{\lambda \in L} X_{\lambda}\) that is the product of the
\end{enumerate}

measures \(\mu_{\lambda}\) for \(\lambda \in J\) and the measures \(\varepsilon_{a_{\lambda}}\) for \(\lambda \in L - J\) . Show that, with respect to the directed set of finite subsets of \(L\) , the measure \(\nu(J)\) tends vaguely to \(\mu = \bigotimes_{\lambda \in L} \mu_{\lambda}\)

but does not tend strongly to \(\mu\) .

¶ 6) With the notations of No.~6 show that, for there to exist a measure \(\mu \neq 0\) on X such that

\[
\mu \left(f _ {J} \circ \operatorname{pr} _ {J}\right) = \lim _ {K \supset J} \mu_ {K} \left(f _ {J} \circ \operatorname{pr} _ {J, K}\right) = \mu_ {J} \left(f _ {J}\right) \cdot \prod_ {\lambda \in L - J} \mu_ {\lambda} (1)\tag{1}
\]

for every finite subset \(\mathbf{J}\) of \(\mathbf{L}\) and every function \(f_{\mathbf{J}} \in \mathcal{C}(\mathrm{X}_{\mathbf{J}}; \mathbf{C})\) , it is necessary and sufficient that the following three conditions be satisfied:

\(1^{\circ} \mu_{\lambda} \neq 0\) for all \(\lambda \in L\) .

\(2^{\circ}\) There exists a finite subset \(\mathbf{J}_0\) of \(\mathbf{L}\) such that the family \((\mu_{\lambda}(1))_{\lambda \in \mathbf{L} - \mathbf{J}_0}\) is multipliable in \(\mathbf{C}^*\) (hence has a product \(\neq 0\) ).

\(3^{\circ}\) The family \((\| \mu_{\lambda}\|)_{\lambda \in L}\) is multipliable in \(\mathbf{R}_{+}^{*}\) (hence has a product \(\neq 0\) ).

Under what conditions is the second member of (1) equal to 0 for every finite subset J of L and every function \(f_{\mathrm{J}} \in \mathcal{C}(\mathrm{X}_{\mathrm{J}}; \mathbf{C})\) ?

\begin{enumerate}
\def\labelenumi{\arabic{enumi})}
\setcounter{enumi}{6}
\tightlist
\item
  With the notations of Exercise 6, assume that the conditions of the exercise are satisfied. Show that for every function \(f \in \mathcal{C}(\mathrm{X};\mathbf{C})\) ,
\end{enumerate}

\[
\int f d \mu = \lim _ {\mathrm{J}} \int f (x _ {\mathrm{J}}, x _ {\mathrm{L} - \mathrm{J}}) d \mu_ {\mathrm{J}} (x _ {\mathrm{J}}),
\]

where the limit is taken with respect to the directed set of finite subsets of L, and where \(x_{J}\) denotes \(\Pr_{J}(x)\) for every \(x \in X\) and every \(J \subset L\) ; similarly,

\[
f (x) = \lim _ {\mathrm{J}} \int f (x _ {\mathrm{J}}, x _ {\mathrm{L} - \mathrm{J}}) d \mu_ {\mathrm{L} - \mathrm{J}} (x _ {\mathrm{L} - \mathrm{J}})
\]

for every \(x \in X\) , provided that all of the measures \(\mu_{\lambda} (\lambda \in L)\) have total mass 1 (where, for every subset \(K\) of \(L\) , \(\mu_{K}\) denotes the product of the subfamily of measures \((\mu_{\lambda})_{\lambda \in K}\) ).

\begin{enumerate}
\def\labelenumi{\arabic{enumi})}
\setcounter{enumi}{7}
\tightlist
\item
  With the notations of Exercise 6, show that if the product of the family \((\mu_{\lambda})\) is defined and \(\neq 0\) , then there exists a countable subset \(K\) of \(L\) such that for every \(\lambda \in L - K\) , the measure \(\mu_{\lambda}\) is positive and of total mass 1 (make use of Exer. 18 of §1).
\end{enumerate}

